% !TeX root = ../main.tex
% 第 3 章 整数维数集合的结构 (对应原书 pp. 28--53)

\chapter{整数维数集合的结构}\label{chap:3}

\section{引言}

本章讨论 $s$ 为整数时 $\mathbb{R}^n$ 中 $s$-集的密度与切线结构.由\cref{cor:2.10} 知 $s$-集分裂为一个正则部分与一个非正则部分，而且我们发现这两类集合表现出显著不同的性质.我们的目标之一是把正则集刻画为可求长曲线或曲面的可数并的子集，从而把测度论的观念与描述性的拓扑观念联系起来.

我们详细叙述 $\mathbb{R}^2$ 中线性可测集（即 $1$-集）的理论.这一工作几乎完全归功于 Besicovitch（\citeyear{besicovitch1928a,besicovitch1938}），后一篇文章包含了一些改进的证明以及进一步的结果.Besicovitch 的证明大多难以超越——除了相对次要的改进，以及在表述上或许可以改进之处.确实，Besicovitch 所用的某些几何方法精巧到令人惊讶：它们当初居然能被想出来.\textcite{deguzman1981} 也描述了本章的部分工作.

\section{曲线与连续统}

正则 $1$-集与可求长曲线紧密相关.事实上，除去一个测度为零的集合不计，正则 $1$-集是可求长曲线的可数族的子集.本节致力于研究曲线，主要从拓扑的观点出发，并与具有有限线性测度的连续统相联系.我们在 $\mathbb{R}^n$ 中工作，因为这一理论并不比平面曲线的情形更复杂.

\newterm{curve}{曲线}{curve}（或 \newterm{jordan-curve}{Jordan 曲线}{Jordan curve}）$\Gamma$ 是连续单射 $\psi:[a,b]\to\mathbb{R}^n$ 的像，其中 $[a,b]\subset\mathbb{R}$ 为闭区间.任何曲线都是\newterm{continuum}{连续统}{continua}，即紧连通集.这是因为紧连通集的连续像既紧又连通.特别地，任何曲线都是 Borel 集，从而是 $\mathcal{H}^s$-可测的.此外，紧集之间的连续双射有连续的逆，故曲线也可以定义为闭区间在同胚意义下的像.

曲线 $\Gamma$ 的\newterm{length}{长度}{length}定义为
\begin{equation}\label{eq:3.1}
  \mathscr{L}(\Gamma)=\sup\sum_{i=1}^{m}|\psi(t_i)-\psi(t_{i-1})|,
\end{equation}
其中上确界取遍 $[a,b]$ 的一切分割 $a=t_0<t_1<\cdots<t_m=b$.若 $\mathscr{L}(\Gamma)<\infty$（即 $\psi$ 具有\newterm{bounded-variation}{有界变差}{bounded variation}），则称 $\Gamma$ 是\newterm{rectifiable}{可求长的}{rectifiable curve}.

注意我们的定义排除了\newterm{curve-self-intersecting}{自交曲线}{curve, self-intersecting}；下面的引理把它们涵盖在内.

\begin{lemma}\label{lem:3.1}
设 $\psi:[a,b]\to\mathbb{R}^n$ 为连续映射，$\psi(a)\neq\psi(b)$.则 $\psi[a,b]$ 包含一条连接 $\psi(a)$ 与 $\psi(b)$ 的曲线.
\end{lemma}

\begin{proof}
对 $[a,b]$ 的每个重点 $x$，令 $I_x$ 是满足 $\psi(t_1)=\psi(t_2)=x$ 的最大（闭）区间 $[t_1,t_2]$.令 $\mathcal{I}$ 表示不被其他此类区间包含的那些区间之族.则 $\mathcal{I}$ 由可数多个两两不交的真闭区间组成.于是可以构造连续满射 $f:[a,b]\to[0,1]$，使 $f(a)=0$，$f(b)=1$，并且当 $t_1\leq t_2$ 时 $f(t_1)\leq f(t_2)$，其中等号成立当且仅当 $t_1$ 与 $t_2$ 位于 $\mathcal{I}$ 的同一区间中.定义 $\psi_0:[0,1]\to\mathbb{R}^n$ 如下：若对某个 $I_x\in\mathcal{I}$ 有 $f^{-1}(u)=I_x$，则令 $\psi_0(u)=x$；否则令 $\psi_0(u)=\psi(f^{-1}(u))$.容易验证 $\psi_0$ 是连续单射，且 $\psi_0(0)=\psi(a)$，$\psi_0(1)=\psi(b)$.
\end{proof}

\eqref{eq:3.1} 中的和不会因添加分割点而减小.显然，若把 $\psi[a,b]$ 分成两条曲线 $\psi[a,c]$ 与 $\psi[c,b]$（$a<c<b$），则新曲线长度之和等于原曲线的长度，即 $\mathscr{L}$ 在有公共端点的曲线上是可加的.

\newterm{curve-arc-length}{用弧长参数化可求长曲线}{curve, parameterized by arc-length}总是可能的，即把 $\Gamma$ 表示为函数 $\psi_0:[0,\mathscr{L}(\Gamma)]\to\mathbb{R}^n$ 的像，使 $\psi_0[0,t]$ 的长度为 $t$.这可以由令 $\psi_0(t)$ 为满足 $\mathscr{L}(\psi[a,u])=t$ 的唯一一点 $\psi(u)$ 而实现.若 $\psi$ 用弧长表示可求长曲线 $\Gamma$，则由 \eqref{eq:3.1} 可见
\begin{equation}\label{eq:3.2}
  |\psi(t_1)-\psi(t_2)|\leq|t_1-t_2|.
\end{equation}
特别地，这蕴含 $\psi$ 是绝对连续函数.

关于曲线定义的进一步讨论见 \textcite{burkillburkill1970}（第 246 页）或 \textcite{pelling1977}.

\begin{lemma}\label{lem:3.2}
若 $\Gamma$ 是曲线，则 $\mathcal{H}^1(\Gamma)=\mathscr{L}(\Gamma)$.
\end{lemma}

\begin{proof}
设 $\Gamma$ 是连接 $z$ 与 $w$ 的曲线.若 $\operatorname{proj}$ 表示从 $\mathbb{R}^n$ 到过 $z$ 与 $w$ 的直线上的正交投影，则当 $x,y\in\mathbb{R}^n$ 时 $|\operatorname{proj}x-\operatorname{proj}y|\leq|x-y|$.由\cref{lem:1.8}，
\[
  \mathcal{H}^1(\Gamma)\geq\mathcal{H}^1(\operatorname{proj}\Gamma)\geq\mathcal{H}^1([z,w])=\mathscr{L}^1([z,w])=|z-w|,
\]
因为 $\operatorname{proj}\Gamma\supset[z,w]$.

现在设 $\Gamma$ 由 $\psi:[a,b]\to\mathbb{R}^n$ 定义.由上面的注记，对任意 $t$ 与 $u$ 有 $\mathcal{H}^1(\psi[t,u])\geq|\psi(t)-\psi(u)|$.于是若 $a=t_0<t_1<\cdots<t_m=b$ 是 $[a,b]$ 的任一分割，则
\[
  \sum_{i}|\psi(t_i)-\psi(t_{i-1})|\leq\sum_{i}\mathcal{H}^1\bigl(\psi[t_{i-1},t_i]\bigr)=\mathcal{H}^1(\Gamma),
\]
因为 $\Gamma$ 的弧 $\psi[t_{i-1},t_i]$ 除端点外互不相交.故 $\mathscr{L}(\Gamma)\leq\mathcal{H}^1(\Gamma)$.

最后，设 $\mathscr{L}(\Gamma)<\infty$，并设 $\psi$ 用弧长参数化 $\Gamma$.由于 $\psi$ 是从 $[0,\mathscr{L}(\Gamma)]$ 到 $\Gamma$ 的满足 \eqref{eq:3.2} 的满射，\cref{lem:1.8} 蕴含
$\mathcal{H}^1(\Gamma)\leq\mathcal{H}^1([0,\mathscr{L}(\Gamma)])=\mathscr{L}(\Gamma)$.
\end{proof}

这条引理的一个推论是：若 $\Gamma$ 是可求长曲线，则当 $s<1$ 时 $\mathcal{H}^s(\Gamma)$ 为无穷，当 $s>1$ 时为零（见 \eqref{eq:1.14}）.

\begin{corollary}\label{cor:3.3}
设 $\psi$ 用弧长参数化可求长曲线 $\Gamma$.若 $E$ 是 $[0,\mathscr{L}(\Gamma)]$ 的任意勒贝格可测子集，则 $\psi(E)$ 是 $\mathbb{R}^n$ 的 $\mathcal{H}^1$-可测子集，且 $\mathcal{H}^1(\psi(E))=\mathscr{L}^1(E)$.
\end{corollary}

\begin{proof}
由\cref{lem:3.2}，可加集函数 $\mathcal{H}^1(\psi(\cdot))$ 与 $\mathscr{L}^1(\cdot)$ 在闭区间上一致，故由通常的测度扩张过程（例如见 \textcite{kingmantaylor1966}），它们在 $[0,\mathscr{L}(\Gamma)]$ 的勒贝格可测子集上一致.
\end{proof}

这条推论使得直线上的勒贝格测度的所有常用结果都可以移植到曲线上.例如，由勒贝格密度定理（\cref{thm:1.13}）可得：若 $E$ 是曲线 $\Gamma$ 的 $\mathcal{H}^1$-可测子集，则在几乎所有的 $x\in E$ 处，当 $\mathcal{H}^1(I)\to 0$ 时有 $\mathcal{H}^1(E\cap I)/\mathcal{H}^1(I)\to 1$，其中 $I$ 是 $\Gamma$ 中包含 $x$ 的子弧.

我们不再进一步使用这条推论.

我们经常要用到下一条引理：

\begin{lemma}\label{lem:3.4}
设 $E$ 是包含 $x$ 与 $y$ 的连续统.若 $|x-y|=\rho$，则 $\mathcal{H}^1\bigl(E\cap B_\rho(x)\bigr)\geq\rho$.特别地，$\mathcal{H}^1(E)\geq|E|$.
\end{lemma}

\begin{proof}
令 $f:\mathbb{R}^n\to[0,\infty)$ 定义为 $f(z)=|z-x|$.则 $f$ 是满足
\begin{equation}\label{eq:3.3}
  |f(z)-f(w)|\leq|z-w|\qquad(z,w\in\mathbb{R}^n)
\end{equation}
的连续映射.集合 $f\bigl(E\cap B_\rho(x)\bigr)$ 包含区间 $[0,\rho]$；否则，对某个 $r$（$0<r<\rho$），$E=(E\cap B_r(x))\cup(E\setminus B_r(x))$ 便成为 $E$ 的由两个不交闭集构成的分解.对映射 $f$ 应用\cref{lem:1.8}，
\[
  \mathcal{H}^{1}\bigl(E\cap B_{\rho}(x)\bigr)
  \geq\mathcal{H}^{1}\bigl(f(E\cap B_{\rho}(x))\bigr)
  \geq\mathcal{H}^{1}([0,\rho])=\mathscr{L}^{1}([0,\rho])=\rho. \qedhere
\]
\end{proof}

现在可以推出可求长曲线的基本密度性质.

\begin{lemma}\label{lem:3.5}
可求长曲线是正则 $1$-集.
\end{lemma}

\begin{proof}
若 $\Gamma$ 可求长，则 $\mathscr{L}(\Gamma)<\infty$，故由\cref{lem:3.2}，$\mathcal{H}^1(\Gamma)<\infty$.由于 $\Gamma$ 至少含两个不同的点，\cref{lem:3.4} 蕴含 $\mathcal{H}^1(\Gamma)>0$，故 $\Gamma$ 是一个 $1$-集.

设 $x$ 是 $\Gamma$ 上非端点的一点，并设 $x$ 把 $\Gamma$ 分成两条可求长子曲线 $\Gamma_-$ 与 $\Gamma_+$.当 $\rho$ 充分小时，由\cref{lem:3.4} 有 $\mathcal{H}^1\bigl(\Gamma_-\cap B_\rho(x)\bigr)\geq\rho$ 与 $\mathcal{H}^1\bigl(\Gamma_+\cap B_\rho(x)\bigr)\geq\rho$，从而 $\mathcal{H}^1\bigl(\Gamma\cap B_\rho(x)\bigr)\geq 2\rho$.于是在 $\Gamma$ 上除端点外的一切点处 $D^1(\Gamma,x)\geq 1$.结合\cref{cor:2.5}，这说明在几乎所有的 $x\in\Gamma$ 处 $D^1(\Gamma,x)$ 存在且等于 $1$.
\end{proof}

接下来讨论可求长曲线的切线的存在性.称 $\mathbb{R}^n$ 中的 $s$-集 $E$ 在 $x$ 处沿方向 $\pm\theta$ \newterm{tangent}{有切线}{tangent}，如果 $\bar{D}^s(E,x)>0$，而且对每个角 $\phi>0$，
\begin{equation}\label{eq:3.4}
  \lim_{r\to 0}r^{-s}\mathcal{H}^{s}\bigl(E\cap(B_r(x)\setminus S_r(x,\theta,\phi)\setminus S_r(x,-\theta,\phi))\bigr)=0
\end{equation}
（沿方向 $\pm\theta$ 过 $x$ 的直线当然就是切线）.显然，$s$-集在它的每一点处至多有一条切线.

我们给出 \textcite{besicovitch1944} 关于可求长曲线在几乎所有点处有切线的初等证明.另一种途径是诉诸有界变差函数可微性的标准结果，见 \textcite{kingmantaylor1966}（9.1~节）.

\begin{lemma}\label{lem:3.6}
设 $\Gamma$ 是端点为 $x$ 与 $y$ 的可求长曲线，$\phi$ 为正角.设 $E$ 是 $\Gamma$ 上这样的点的集合：$x$ 属于 $\Gamma$ 的任意小的一对子弧，这对子弧所对的弦彼此的夹角大于 $2\phi$.则 $\mathcal{H}^1(E)\leq\bigl(\mathscr{L}(\Gamma)-|x-y|\bigr)/(1-\cos\phi)$.
\end{lemma}

\begin{proof}
设 $L$ 表示过 $x$ 与 $y$ 的直线，并设 $\mathcal{V}$ 为 $\Gamma$ 的这样的闭子弧之族：子弧所对的弦与 $L$ 的夹角大于 $\phi$.由引理的条件，$\mathcal{V}$ 是 $E$ 的 Vitali 类.于是利用覆盖定理（\cref{thm:1.10}(b)），对任意 $\varepsilon>0$，可找到 $\Gamma$ 的属于 $\mathcal{V}$ 的不交子弧的有限族 $\Gamma_1,\ldots,\Gamma_m$ 使
\begin{equation}\label{eq:3.5}
  \mathcal{H}^{1}(E)\leq\sum_{i=1}^{m}|\Gamma_i|+\varepsilon\leq\sum_{i=1}^{m}\mathscr{L}(\Gamma_i)+\varepsilon.
\end{equation}
若 $\Gamma_0',\Gamma_1',\ldots,\Gamma_m'$ 是（可能退化为空的）补弧，则正交投影到 $L$ 上并利用投影不增长度的平凡事实，可见
\[
  \cos\phi\sum_{i}\mathscr{L}(\Gamma_i)+\sum_{i}\mathscr{L}(\Gamma_i')\geq|x-y|.
\]
由于 $\sum_{i}\mathscr{L}(\Gamma_i)+\sum_{i}\mathscr{L}(\Gamma_i')=\mathscr{L}(\Gamma)$，用 \eqref{eq:3.5} 得
\[
  \mathcal{H}^{1}(E)-\varepsilon\leq\sum_{i}\mathscr{L}(\Gamma_i)\leq\bigl(\mathscr{L}(\Gamma)-|x-y|\bigr)/(1-\cos\phi).
\]
$\varepsilon$ 任意，证明完成.
\end{proof}

\begin{corollary}\label{cor:3.7}
若 $\phi>0$，$E$ 是可求长曲线 $\Gamma$ 上这样的点的集合：$x$ 属于 $\Gamma$ 的任意小的一对子弧，这对子弧所对的弦彼此的夹角大于 $2\phi$，则 $\mathcal{H}^1(E)=0$.
\end{corollary}

\begin{proof}
给定 $\varepsilon>0$，由 \eqref{eq:3.1} 可在 $\Gamma$ 上（按次序，且 $x_0$ 与 $x_m$ 为 $\Gamma$ 的端点）找到点 $x_0,x_1,\ldots,x_m$ 使
\[
  \mathscr{L}(\Gamma)<\sum_{i=1}^{m}|x_i-x_{i-1}|+\varepsilon.
\]
记 $\Gamma_i$ 为 $\Gamma$ 介于 $x_{i-1}$ 与 $x_i$ 之间的部分.依次对每个 $\Gamma_i$ 应用\cref{lem:3.6} 得
\begin{align*}
  \mathcal{H}^{1}(E)
  &=\sum_{i=1}^{m}\mathcal{H}^{1}(E\cap\Gamma_i)
  \leq\sum_{i=1}^{m}\bigl(\mathscr{L}(\Gamma_i)-|x_i-x_{i-1}|\bigr)/(1-\cos\phi)\\
  &=\bigl(\mathscr{L}(\Gamma)-\sum_{i=1}^{m}|x_i-x_{i-1}|\bigr)/(1-\cos\phi)
  <\varepsilon/(1-\cos\phi).
\end{align*}
故 $\mathcal{H}^1(E)=0$.
\end{proof}

\begin{theorem}\label{thm:3.8}
可求长曲线 $\Gamma$ 在它的几乎所有点处有切线.
\end{theorem}

\begin{proof}
由于连续统 $\Gamma$ 至少有两个点，\cref{lem:3.4} 蕴含在所有 $x\in\Gamma$ 处 $\bar{D}^s(E,x)\geq\frac{1}{2}$.

设 $\psi:[a,b]\to\mathbb{R}^n$ 是 $\Gamma$ 的一个定义函数.由\cref{cor:3.7}，对几乎所有的 $x\in\Gamma$ 可找到单位向量 $\theta$，使得给定 $\phi>0$，当 $|u-t|<\varepsilon$ 时（这里 $\psi(t)=x$）有
\[
  \psi(u)\in S(x,\theta,\phi)\cup S(x,-\theta,\phi).
\]
此外，若 $\rho$ 为正数，则当 $|u-t|\geq\varepsilon$ 时 $\psi(u)\notin B_\rho(x)$（否则可选取满足 $|u_i-t|\geq\varepsilon$ 且 $\psi(u_i)\to x$ 的数列 $\{u_i\}$；利用序列紧性，存在 $u$ 使 $|u-t|\geq\varepsilon$ 且 $\psi(u)=x=\psi(t)$，这就给出了 $\Gamma$ 的一个重点）.于是
\[
  \Gamma\cap\bigl(B_{\rho}(x)\setminus S_{\rho}(x,\theta,\phi)\setminus S_{\rho}(x,-\theta,\phi)\bigr)=\emptyset,
\]
从而 \eqref{eq:3.4} 作为平凡推论成立.
\end{proof}

曲线的切线常用可微性来理解.\cref{thm:3.8} 证明的一个微小变形表明：由函数 $\psi$ 定义的可求长曲线在 $\psi(t)$ 处有切线，只要导数 $\psi'(t)$ 存在且不为零.

Besicovitch（\citeyear{besicovitch1934a,besicovitch1956,besicovitch1957,besicovitch1960}） 还考虑了如何把切线的定义调整后用于具有无穷线性测度的曲线与集合.

含于可求长曲线的可数并中的 $1$-集称为 \newterm{Y-set}{$Y$-集}{$Y$-set}.

\cref{lem:3.5} 与 \cref{thm:3.8} 有针对 $Y$-集的对应物.

\begin{corollary}\label{cor:3.9}
$Y$-集是正则 $1$-集.
\end{corollary}

\begin{proof}
把\cref{cor:2.8} 应用于\cref{lem:3.5}：含于可求长曲线中的 $1$-集几乎处处密度为 $1$.由\cref{cor:2.7}，$Y$-集亦然.
\end{proof}

\begin{corollary}\label{cor:3.10}
$Y$-集在它的几乎所有点处有切线.
\end{corollary}

\begin{proof}
由\cref{cor:3.9}，$Y$-集几乎处处具有正的下密度.把\cref{cor:2.6} 应用于\cref{thm:3.8}：含于可求长曲线中的 $1$-集几乎处处满足 \eqref{eq:3.4}，故由\cref{cor:2.7}，$Y$-集在几乎所有点处有切线.
\end{proof}

最终我们将证明正则 $1$-集是一个 $Y$-集加上一个测度为零的集合，从而由\cref{cor:3.10} 可得任何正则 $1$-集几乎处处有切线.

具有下一条引理所描述性质的集合有时称为 \newterm{lipschitz-set}{Lipschitz 集}{Lipschitz set}.

\begin{lemma}\label{lem:3.11}
设 $E$ 是 $\mathbb{R}^n$ 的有界子集，使得当 $x,y\in E$ 时线段 $[x,y]$ 与固定直线 $L$ 的夹角至多为 $\phi<\frac{1}{2}\pi$.则 $E$ 含于一条可求长曲线之中.
\end{lemma}

\begin{proof}
不妨设 $E$ 是闭集.设 $\Pi(t)$ 为与 $L$ 垂直、距某原点距离为 $t$ 的超平面，并设 $a$ 与 $b$ 是使 $\Pi(t)$ 与 $E$ 相交的 $t$ 的极端值.于是对每个 $t$，$\Pi(t)$ 至多含 $E$ 的一个点；若存在就用 $\psi(t)$ 表示这个点，若不存在且 $a<t<b$，则令 $\psi(t)$ 为 $\Pi(t)$ 上位于连接两侧离 $\Pi(t)$ 最近的 $E$ 中的点的线段上的点.当 $t_1\neq t_2$ 时，线段 $[\psi(t_1),\psi(t_2)]$ 与 $L$ 的夹角至多为 $\phi$，故 $|\psi(t_1)-\psi(t_2)|\leq|t_1-t_2|/\cos\phi$.由 \eqref{eq:3.1}，$\psi$ 定义了一条可求长曲线.
\end{proof}

接下来我们证明，具有有限 $\mathcal{H}^1$-测度的连续统除去一个零测集之外是可求长曲线的可数族，从而是正则的.遵循 \textcite{besicovitch1938}，我们以弧连通作为中间阶段.回忆：称集合 $E$ 是\newterm{arcwise-connected}{弧连通的}{arcwise connected}，如果任给 $z,w\in E$，存在位于 $E$ 内连接 $z$ 与 $w$ 的曲线.紧的弧连通集显然是连续统.我们需要下述逆命题.

\begin{lemma}\label{lem:3.12}
满足 $\mathcal{H}^1(E)<\infty$ 的连续统 $E$ 是弧连通的.
\end{lemma}

\begin{proof}
这本质上是\cref{lem:3.4} 的推论.取 $z,w\in E$.由于 $E$ 连通，$E$ 是\newterm{chain-connected}{链连通的}{chain connected}，即对每个 $\varepsilon>0$ 可在 $E$ 中找到点链 $z=x_0,x_1,\ldots,x_m=w$ 使 $|x_i-x_{i-1}|\leq\varepsilon$（$1\leq i\leq m$）.（这可以通过证明与 $z$ 链连通的点之集在 $E$ 中既开又闭而得.）必要时删去链的中间点，可设当 $|i-j|\geq 2$ 时 $|x_i-x_j|>\varepsilon$.于是 $\mathbb{R}^n$ 中没有点落在多于两个球 $B_{\frac{1}{2}\varepsilon}(x_i)$ 之中，故设 $m\geq 2$ 并用\cref{lem:3.4} 得
\[
  2\mathcal{H}^{1}(E)\geq\sum_{i=0}^{m}\mathcal{H}^{1}\bigl(E\cap B_{\frac{1}{2}\varepsilon}(x_i)\bigr)\geq\frac{1}{2}m\varepsilon.
\]
于是，若 $\Gamma_\varepsilon$ 表示连接 $x_0,x_1,\ldots,x_m$ 所得的折线曲线（必然不自交），则
\begin{equation}\label{eq:3.6}
  \mathscr{L}(\Gamma_{\varepsilon})=\sum_{i=1}^{m}|x_i-x_{i-1}|\leq m\varepsilon\leq 4\mathcal{H}^{1}(E).
\end{equation}
对 $0\leq t\leq 1$，令 $\psi_\varepsilon(t)$ 为 $\Gamma_\varepsilon$ 上沿其长度比例 $t$ 处的点，使 $\Gamma_\varepsilon$ 介于 $z=\psi_\varepsilon(0)$ 与 $\psi_\varepsilon(t)$ 之间的部分长度为 $t\mathscr{L}(\Gamma_\varepsilon)$.于是当 $0\leq t_1<t_2\leq 1$ 时，由 \eqref{eq:3.6}，
\begin{equation}\label{eq:3.7}
  |\psi_{\varepsilon}(t_1)-\psi_{\varepsilon}(t_2)|\leq\mathscr{L}\bigl(\psi_{\varepsilon}[t_1,t_2]\bigr)\leq|t_2-t_1|\,4\mathcal{H}^{1}(E).
\end{equation}
曲线族 $\{\Gamma_\varepsilon\}_{\varepsilon\leq 1}$ 都位于 $\mathbb{R}^n$ 的一个有界子集之中，故由 $[0,1]$ 到 $\mathbb{R}^n$ 的函数族 $\{\psi_\varepsilon\}_{\varepsilon\leq 1}$ 一致有界，且由 \eqref{eq:3.7} 等度连续.由 Arzelà--Ascoli 定理（例如见 \textcite{dunfordschwartz1958}，第四部分 6.7 节），这样的函数族是序列紧的，即存在序列 $\varepsilon(j)\to 0$ 与连续函数 $\psi:[0,1]\to\mathbb{R}^n$ 使 $\psi_{\varepsilon(j)}$ 在 $[0,1]$ 上一致收敛于 $\psi$.令 $\psi[0,1]=\Gamma$.显然 $\psi(0)=z$，$\psi(1)=w$.又若 $x\in\Gamma$，则给定 $\delta>0$，可找到 $\varepsilon(j)<\frac{1}{2}\delta$ 使 $\Gamma_{\varepsilon(j)}$ 与球 $B_{\frac{1}{2}\delta}(x)$ 相交.于是 $B_\delta(x)$ 中有 $E$ 的某个点（即 $\Gamma_{\varepsilon(j)}$ 的一个顶点）.由于 $E$ 是闭的，我们得出 $x\in E$，故 $\Gamma\subset E$.若 $\psi$ 是单射，则 $\Gamma$ 是曲线.否则，由\cref{lem:3.1} 可把 $\Gamma$ 缩减为连接 $z$ 与 $w$ 的一条曲线.
\end{proof}

\begin{lemma}\label{lem:3.13}
任何满足 $\mathcal{H}^1(E)<\infty$ 的紧弧连通集 $E$ 由可求长曲线的可数并再加上一个 $\mathcal{H}^1$-测度为零的集合组成.
\end{lemma}

\begin{proof}
归纳地定义曲线序列 $\{\Gamma_j\}$.令 $\Gamma_1$ 为 $E$ 中连接 $E$ 的两个最远点的曲线.设曲线 $\Gamma_1,\Gamma_2,\ldots,\Gamma_k$ 已经定义.令 $x$ 为 $E$ 中距 $\bigcup_{1}^{k}\Gamma_j$ 最远（距离记为 $d_k$）的点（由于所涉集合是紧的，该最大值可以取得）.若 $d_k=0$，过程终止，结论得证.否则，令 $\Gamma_{k+1}$ 为 $E$ 中连接 $x$ 与 $\bigcup_{1}^{k}\Gamma_j$ 的曲线，且除一个端点外 $\Gamma_{k+1}$ 与 $\bigcup_{1}^{k}\Gamma_j$ 不交.由\cref{lem:3.4},
\begin{equation}\label{eq:3.8}
  \mathcal{H}^{1}(\Gamma_{k+1})\geq d_k.
\end{equation}
于是
\begin{equation}\label{eq:3.9}
  \sum_{j=1}^{\infty}d_j\leq\sum_{1}^{\infty}\mathcal{H}^{1}(\Gamma_j)\leq\mathcal{H}^{1}(E)<\infty,
\end{equation}
故 $d_j\to 0$.若 $x\in E$ 且 $x$ 到 $\bigcup_{1}^{\infty}\Gamma_j$ 的距离 $d$ 为正，则对某个 $j$ 有 $d_j<d$，而 $x$ 本应被选为 $\Gamma_{j+1}$ 的端点，矛盾.故 $E$ 是 $\bigcup_{1}^{\infty}\Gamma_j$ 的闭包.由\cref{lem:3.2}, $\mathscr{L}(\Gamma_j)=\mathcal{H}^1(\Gamma_j)\leq\mathcal{H}^1(E)<\infty$，故每个 $\Gamma_j$ 都是可求长曲线.最后，对每个 $k$，$\bigcup_{1}^{k}\Gamma_j$ 是闭集，故
\[
  \mathcal{V}=\Bigl\{B_r(x):x\in E\setminus\bigcup_{1}^{k}\Gamma_j
  \ \text{且}\ B_r(x)\cap\bigcup_{1}^{k}\Gamma_j=\emptyset\Bigr\}
\]
是 $E\setminus\bigcup_{1}^{k}\Gamma_j$ 的由球构成的 Vitali 类.设 $x$ 为某个 $B\in\mathcal{V}$ 的圆心.由于 $x$ 位于 $\bigcup_{k+1}^{\infty}\Gamma_j$ 的闭包中，有该并集的点任意接近 $x$.这些点必须经由 $\bigcup_{k+1}^{\infty}\Gamma_j$ 中的一列弧连接到 $\bigcup_{1}^{k}\Gamma_j$ 中必在 $B$ 之外的一点，故由\cref{lem:3.4},
\[
  \frac{1}{2}|B|\leq\mathcal{H}^{1}\Bigl(B\cap\bigcup_{k+1}^{\infty}\Gamma_j\Bigr)
\]
对任何 $B\in\mathcal{V}$ 成立.于是给定 $\varepsilon>0$，可用 Vitali 定理（\cref{thm:1.10}(b)）从 $\mathcal{V}$ 中选出不交球列 $\{B_i\}$ 使
\[
  \mathcal{H}^{1}\Bigl(E\setminus\bigcup_{1}^{k}\Gamma_j\Bigr)
  \leq\sum_{i}|B_i|+\varepsilon
  \leq 2\sum_{i}\mathcal{H}^{1}\Bigl(B_i\cap\bigcup_{k+1}^{\infty}\Gamma_j\Bigr)+\varepsilon
  \leq 2\mathcal{H}^{1}\Bigl(\bigcup_{k+1}^{\infty}\Gamma_j\Bigr)+\varepsilon.
\]
从而
\[
  \mathcal{H}^{1}\Bigl(E\setminus\bigcup_{1}^{k}\Gamma_j\Bigr)\leq 2\mathcal{H}^{1}\Bigl(\bigcup_{k+1}^{\infty}\Gamma_j\Bigr).
\]
由 \eqref{eq:3.9}，当 $k\to\infty$ 时右端趋于零，故 $\mathcal{H}^1\bigl(E\setminus\bigcup_{1}^{\infty}\Gamma_j\bigr)=0$，如所需.
\end{proof}

上面的引理看起来已由 \textcite{wazewski1927} 独立于 Besicovitch 证明.

\begin{theorem}\label{thm:3.14}
设 $E$ 是满足 $\mathcal{H}^1(E)<\infty$ 的连续统.则 $E$ 由可求长曲线的可数并再加上一个 $\mathcal{H}^1$-测度为零的集合组成.
\end{theorem}

\begin{proof}
合并\cref{lem:3.12,lem:3.13}.
\end{proof}

\begin{corollary}\label{cor:3.15}
设 $E$ 是满足 $\mathcal{H}^1(E)<\infty$ 的连续统.则 $E$ 是正则的，且在几乎所有点处有切线.
\end{corollary}

\begin{proof}
由\cref{thm:3.14}, $E$ 是一个 $Y$-集加上一个零测集，故结论由\cref{cor:3.9,cor:3.10} 立即得到.（零测集不影响密度与角密度的值.）
\end{proof}

人们还建立了若干其他关于把正则 $1$-集嵌入曲线的结果.例如，\cref{thm:3.14} 中的例外点集可以视为一条可能自交的可求长曲线的长度为 $0$ 的子集.类似地，任何有限测度的连续统 $E$ 可嵌入一条长度至多为 $2\mathcal{H}^1(E)$ 的（可自交的）可求长曲线（见\cref{exe:3.5}）.

本节最后证明一个最初归功于 Gołąb 的关于连续统线性测度\newterm{golab-semicontinuity}{半连续性的定理}{semicontinuity theorem, Gołąb's}.这一结果在几何测度论中相当重要，在几何问题中它被用来断言具有指定性质的测度最大的连续统或曲线的存在性.我们采用 \textcite{blaschke1916} 著名的选择定理；\textcite{rogers1970} 在更广泛的框架下讨论了该定理.

若 $E\subset\mathbb{R}^n$，则 $E$ 的 $\delta$-\newterm{delta-parallel-body}{平行体}{delta-parallel body}是距 $E$ 的距离不超过 $\delta$ 的点构成的闭集，即
\[
  [E]_{\delta}=\Bigl\{x\in\mathbb{R}^n:\inf_{y\in E}|x-y|\leq\delta\Bigr\}.
\]
\newterm{hausdorff-metric}{Hausdorff 度量}{Hausdorff metric} $\delta$ 定义在 $\mathbb{R}^n$ 的所有非空紧子集之族上：
\[
  \delta(E,F)=\inf\bigl\{\delta:E\subset[F]_{\delta}\ \text{且}\ F\subset[E]_{\delta}\bigr\}.
\]
验证 $\delta$ 是一个度量是简单的练习.

\begin{theorem}[\newterm{blaschke-selection-theorem}{Blaschke 选择定理}{Blaschke selection theorem}]\label{thm:3.16}
设 $\mathcal{C}$ 是位于 $\mathbb{R}^n$ 的有界部分 $B$ 中的非空紧集的一个无穷族.则存在 $\mathcal{C}$ 中互不相同的集合的序列 $\{E_j\}$，在 Hausdorff 度量下收敛于一个非空紧集 $E$.
\end{theorem}

\begin{proof}
首先从 $\mathcal{C}$ 中产生一个 Cauchy 序列.设 $\{E_{1,i}\}_i$ 是 $\mathcal{C}$ 中互不相同的集合的任意序列.对每个 $k>1$，如下定义 $\{E_{k-1,i}\}_i$ 的无穷子列 $\{E_{k,i}\}_i$：令 $\mathcal{B}_k$ 为覆盖 $B$ 的直径至多为 $1/k$ 的闭球的有限族.$E_{k-1,i}$ 中的每一个都与这些球的某个特定组合相交，故由抽屉原理（不同的组合只有有限多种），必存在 $\{E_{k-1,i}\}_i$ 的无穷子族 $\{E_{k,i}\}_i$，它们都恰好与 $\mathcal{B}_k$ 中同样的球相交.若 $F$ 是 $\mathcal{B}_k$ 中这一特定组合的球之并，则对一切 $i$ 有 $E_{k,i}\subset F\subset[E_{k,i}]_{1/k}$，于是 $\delta(E_{k,i},F)\leq 1/k$，从而对一切 $i,j$ 有 $\delta(E_{k,i},E_{k,j})\leq 2/k$.令 $E_i=E_{i,i}$，立即看到
\begin{equation}\label{eq:3.10}
  \delta(E_i,E_j)\leq 2/\min\{i,j\},
\end{equation}
故 $\{E_i\}_i$ 是 Cauchy 序列.

现在
\[
  E=\bigcap_{j=1}^{\infty}\overline{\bigcup_{i=j}^{\infty}E_i}
\]
是非空紧集，因为它是非空紧集的递减序列之交.（横线表示取闭包.）由 \eqref{eq:3.10}, $\overline{\bigcup_{i=j}^{\infty}E_i}\subset[E_j]_{2/j}$，故对一切 $j$ 有 $E\subset[E_j]_{2/j}$.另一方面，若 $x\in E_j$，则由 \eqref{eq:3.10}，当 $i\geq j$ 时 $x\in[E_i]_{2/j}$，于是当 $k\geq j$ 时 $x\in[\overline{\bigcup_{i=k}^{\infty}E_i}]_{2/j}$.选取 $y_k\in\overline{\bigcup_{i=k}^{\infty}E_i}$ 使 $|x-y_k|\leq 2/j$.由序列紧性，$\{y_k\}$ 的子列收敛于某个 $y\in\mathbb{R}^n$，且 $|x-y|\leq 2/j$.但 $y\in\bigcap_{k=1}^{\infty}\overline{\bigcup_{i=k}^{\infty}E_i}=E$，故 $x\in[E]_{2/j}$.于是 $E_j\subset[E]_{2/j}$，从而 $\delta(E,E_j)\leq 2/j$.故 $\{E_j\}$ 在 Hausdorff 度量下收敛于 $E$.
\end{proof}

注意，上述证明的后半部分本质上就是 $\delta$ 是完备度量的一个证明.

每对点都由唯一的可求长道路连接的连续统称为\newterm{tree}{树}{tree}.下一条引理使我们能够对树的序列证明半连续性定理，并通过逼近把结果推广到一般连续统.

\begin{lemma}\label{lem:3.17}
设 $F\subset\mathbb{R}^n$ 是满足 $\mathcal{H}^1(F)<\infty$ 的树.则给定 $\delta>0$，可写 $F=\bigcup_{i=1}^{k}F_i$，其中各 $F_i$ 是连续统，满足

(a) $\sum_{i=1}^{k}|F_i|\leq\sum_{i=1}^{k}\mathcal{H}^1(F_i)=\mathcal{H}^1(F)$;

(b) 对一切 $i$ 有 $|F_i|\leq\delta$;

(c) $k\leq 3\delta^{-1}\mathcal{H}^1(F)+1$.
\end{lemma}

\begin{proof}
若 $|F|\leq\delta$，可取 $k=1$，$F_1=F$，故设 $|F|>\delta$.设 $x$ 为 $F$ 的任一点，令 $m$ 为 $F$ 中的点到 $x$ 的道路距离的上确界；则由\cref{lem:3.4} 有 $m\leq\mathcal{H}^1(F)<\infty$，而 $m>\frac{1}{2}\delta$，否则将有 $|F|\leq\delta$.设 $y$ 为 $F$ 中与 $x$ 的道路距离大于 $m-\frac{1}{6}\delta$ 的一点，并设 $z$ 为连接 $y$ 与 $x$ 的唯一道路上与 $x$ 的道路距离为 $m-\frac{1}{2}\delta$ 的点.点 $z$ 确定了 $F$ 到两个子树 $F_1$ 与 $F'$ 的一个分割，$z$ 是它们仅有的公共点，其中 $F_1$ 由 $F$ 中到 $x$ 的连接经过 $z$ 的那些点组成，见图~\ref{fig:3.1}. $F_1$ 的每一点与 $z$ 的道路距离在 $\frac{1}{2}\delta$ 之内，故 $|F_1|\leq\delta$.由\cref{lem:3.4}, $|F_1|\leq\mathcal{H}^1(F_1)$，且 $\mathcal{H}^1(F_1)$ 大于 $y$ 到 $z$ 的道路距离，后者至少为 $\bigl(m-\frac{1}{6}\delta\bigr)-\bigl(m-\frac{1}{2}\delta\bigr)=\frac{1}{3}\delta$.

若 $|F'|>\delta$，则对树 $F'$ 重复这一过程以分离出子树 $F_2$，如此继续，直至剩下直径至多为 $\delta$ 的树 $F_k$.结论的 (a) 与 (b) 是直接的，而 (c) 由 (a) 得出，因为当 $1\leq i\leq k-1$ 时 $\frac{1}{3}\delta\leq\mathcal{H}^1(F_i)$.
\end{proof}

\begin{figure}[htbp]
  \centering
  \includegraphics[width=0.5\textwidth]{fig3-1.png}
  \caption{引理 3.17 证明中树 $F$ 的分解}
  \label{fig:3.1}
\end{figure}

\begin{theorem}\label{thm:3.18}
设 $\{E_j\}$ 是 $\mathbb{R}^n$ 中在 Hausdorff 度量下收敛于紧集 $E$ 的连续统序列.则 $E$ 是连续统，且
\[
  \mathcal{H}^{1}(E)\leq\varliminf_{j\to\infty}\mathcal{H}^{1}(E_j).
\]
\end{theorem}

\begin{proof}
若 $E$ 不连通，则有分解 $E=G_1\cup G_2$，其中 $G_1$ 与 $G_2$ 是非空不交闭集，且它们之间的距离严格为正.于是当 $j$ 充分大时，$E_j\subset[G_1]_{\frac{1}{4}\delta}\cup[G_2]_{\frac{1}{4}\delta}$，其中 $[G_1]_{\frac{1}{4}\delta}$ 与 $[G_2]_{\frac{1}{4}\delta}$ 是含有 $E_j$ 的点的不交平行体，这与 $E_j$ 的连通性矛盾.故 $E$ 是连续统.

为证不等式，不妨设对一切 $j$ 有 $\mathcal{H}^1(E_j)\leq c<\infty$.对每个 $j$ 选取 $E_j$ 的有限子集 $S_j$，使 $S_j$ 也在 Hausdorff 度量下收敛于 $E$.（令 $\delta_j\to 0$，取有限集 $S_j$ 使 $E_j$ 的每一点都落在 $S_j$ 中某点的 $\delta_j$ 邻域内.）由\cref{lem:3.12}, 各 $E_j$ 弧连通，故可找到树 $F_j$ 使 $S_j\subset F_j\subset E_j$.（逐次添加``分支''构造 $F_j$，把 $S_j$ 的每一点用 $E_j$ 中的弧连接到已构造的 $F_j$ 部分.）于是在 Hausdorff 度量下 $F_j\to E$.

固定 $\delta>0$.利用\cref{lem:3.17}, 可分解每棵树 $F_j$ 使 $F_j=\bigcup_{i=1}^{k}F_{ji}$，其中对一切 $i,j$ 有 $|F_{ji}|\leq\delta$，且对一切 $j$ 有 $\sum_{i=1}^{k}|F_{ji}|\leq\mathcal{H}^1(F_j)$.（当然可以对每个 $j$ 用同一个满足 $k\leq 3\delta^{-1}c+1$ 的 $k$.）对 $\{F_{ji}\}_j$ 依次对每个 $i$ 应用 Blaschke 选择定理（\cref{thm:3.16}），通过取子列并重新编号，可设 $\{F_{ji}\}_j$ 在 Hausdorff 度量下收敛于非空紧集 $H_i$（$1\leq i\leq k$）.当然对一切 $i$ 有 $|H_i|\leq\delta$，且 $E\subset\bigcup_{1}^{k}H_i$.于是
\[
  \mathcal{H}^{1}_{\delta}(E)\leq\sum_{i}|H_i|
  =\lim_{j\to\infty}\sum_{i}|F_{ji}|
  \leq\varliminf_{j\to\infty}\mathcal{H}^{1}(F_j)
  \leq\varliminf_{j\to\infty}\mathcal{H}^{1}(E_j),
\]
令 $\delta\to 0$ 即得结论.
\end{proof}

下面的选择结果在几何问题中常常有用.


\begin{corollary}\label{cor:3.19}
设 $\mathcal{C}$ 是 $\mathcal{H}^1$-测度至多为 $c$ 且位于 $\mathbb{R}^n$ 的有界部分中的连续统的无穷族.则存在 $\mathcal{C}$ 中互不相同的集合的序列，在 Hausdorff 度量下收敛于满足 $\mathcal{H}^1(E)\leq c$ 的连续统 $E$.
\end{corollary}

\begin{proof}
用\cref{thm:3.16} 从 $\mathcal{C}$ 中得到收敛于非空紧集 $E$ 的序列 $\{E_j\}$.\cref{thm:3.18} 保证 $E$ 是满足 $\mathcal{H}^1(E)\leq c$ 的连续统.
\end{proof}

\section{密度与正则 $1$-集的刻画}

现在把注意力限于平面的子集.本节导出不规则 $1$-集的下密度的一个本质的上界，从而刻画正则与不规则 $1$-集.

若一个 $1$-集与每条可求长曲线的交都具有 $\mathcal{H}^1$-测度零，则称它为 \newterm{Z-set}{$Z$-集}{$Z$-set}.

由\cref{cor:3.9} 可得不规则集是 $Z$-集.我们将证明逆命题也成立.

\begin{lemma}\label{lem:3.20}
$Z$-集与具有有限 $\mathcal{H}^1$-测度的连续统的交测度为零.
\end{lemma}

\begin{proof}
由\cref{thm:3.14} 与 $Z$-集的定义立得.
\end{proof}

本节的主要部分致力于证明 $Z$-集的上密度几乎处处严格小于 $1$，从而是非正则的.为此我们使用 \textcite{besicovitch1938}（第 15 节）的``圆对''思想.（\newterm{circle-pair}{圆对}{circle-pair}是由半径相等的两个圆盘构成的图形，每个圆盘的圆心位于另一个圆盘的圆周上.）证明复杂但相当精妙.我们遵循 \textcite{morserandolph1944}（第 9 节）中给出的 Besicovitch 证明的推广形式.首先需要一个关于从连续统中除去圆盘内部的拓扑结果.回忆：若圆盘族中没有任何成员包含于另一个成员，则称该族是半不交的.

\begin{lemma}\label{lem:3.21}
设 $E$ 是 $\mathbb{R}^2$ 中的连续统.设 $\{B_i\}_{1}^{\infty}$ 是含于 $E$ 中的闭圆盘的可数半不交族，且对任意 $d>0$，满足 $|B_i|\geq d$ 的 $i$ 只有有限多个.则若 $\Gamma_i$ 是 $B_i$ 的圆周，
\[
  F=\Bigl(E\setminus\bigcup_{i}B_i\Bigr)\cup\bigcup_{i}\Gamma_i
\]
是连续统.
\end{lemma}

\begin{proof}
注意 $F=(E\setminus\bigcup_{i}\operatorname{int}B_i)\cup\bigcup_{i}\Gamma_i$，故为证 $F$ 闭，只需证 $\overline{\bigcup_{i}\Gamma_i}\subset F$.设 $x\in\overline{\bigcup_{i}\Gamma_i}$ 而 $x\notin E\setminus\bigcup_{i}\operatorname{int}B_i$，则 $x\in\operatorname{int}B_k$，对某个 $k$.设 $d>0$ 为 $x$ 到 $B_k$ 圆周的距离.若 $\Gamma_j$ 与 $B_{\frac{1}{2}d}(x)$ 相交，则不仅 $B_j$ 与 $B_{\frac{1}{2}d}(x)$ 相交，而且它与 $\mathbb{R}^2\setminus B_k$ 也相交，因为 $B_j\not\subset B_k$.于是 $|\Gamma_j|\geq\frac{1}{2}d$，故只有有限多个圆 $\Gamma_j$ 能与 $B_{\frac{1}{2}d}(x)$ 相交.既然 $x\in\overline{\bigcup_i\Gamma_i}$，我们得出对某个 $k$，$x\in\overline{\bigcup_{i=1}^{k}\Gamma_i}=\bigcup_{i=1}^{k}\Gamma_i\subset F$，如所需.

为证 $F$ 连通，设 $F=F_1\cup F_2$，其中 $F_1,F_2$ 是不交闭集，并由此导出矛盾.令
\[
  E_1=F_1\cup\bigcup_{\{i:\Gamma_i\subset F_1\}}B_i,\qquad
  E_2=F_2\cup\bigcup_{\{i:\Gamma_i\subset F_2\}}B_i.
\]
每个 $\Gamma_i$ 都含于 $E_1$ 或 $E_2$ 之一，故 $E_1\cup E_2=E$.由于圆盘 $\{B_i\}$ 半不交且对每个 $i$ 有 $F\cap B_i\subset\bigcup_{j}\Gamma_j$，容易验证 $E_1\cap E_2=\emptyset$.集合 $E_1$ 是闭的：若 $x\in\overline{\bigcup_{\{i:\Gamma_i\subset F_1\}}B_i}$，则 $x$ 要么是对某个 $k$ 为 $\bigcup_{1}^{k}B_i$ 中点列的极限，要么是圆周落在闭集 $F_1$ 中且半径趋于 $0$ 的一列圆盘的极限点.无论哪种情形都有 $x\in F_1$.类似地 $E_2$ 闭，故 $E=E_1\cup E_2$ 是 $E$ 的由不交闭集构成的分解.由 $E$ 的连通性，$E_1$ 或 $E_2$ 之一为空，从而 $F_1$ 或 $F_2$ 之一为空.
\end{proof}

$Z$-集下密度上界之证明的实质包含在下一条引理中.记 $R(x,y)$ 为以 $x$ 与 $y$ 为心的圆对的\newterm{common-region}{公共区域}{common region}，即 $R(x,y)=\operatorname{int}\bigl(B_{|x-y|}(x)\cap B_{|x-y|}(y)\bigr)$.粗略地说，我们将证明：若以 $1$-集 $E$ 中的点为心的圆对的公共区域包含了 $E$ 的很大一部分，则可以用在很大程度上位于 $E$ 内的曲线把 $E$ 的各个分支连接起来.

\begin{lemma}\label{lem:3.22}
设 $E$ 是 $\mathbb{R}^2$ 中的 $1$-集，$\alpha>0$.设 $E_0$ 是 $E$ 的紧子集，$\mathcal{H}^1(E_0)>0$，使得当 $x_1,x_2\in E_0$ 时 $\mathcal{H}^1\bigl(E\cap R(x_1,x_2)\bigr)\geq\alpha|x_1-x_2|$.则存在连续统 $H$ 使 $0<\mathcal{H}^1(H\cap E)\leq\mathcal{H}^1(H)<\infty$.
\end{lemma}

\begin{proof}
利用通常的一致化过程，由\cref{cor:2.5} 可找到 $\rho_1>0$ 与紧集 $F\subset E_0$，$\mathcal{H}^1(F)>0$，使
\begin{equation}\label{eq:3.11}
  \mathcal{H}^{1}\bigl(B_r(x)\cap E\bigr)\leq 2\cdot 2r
  \quad\text{若}\ x\in F\ \text{且}\ 0<r\leq\rho_1.
\end{equation}
进一步，由\cref{cor:2.4,cor:2.5}，存在点 $y\in F$ 与数 $\rho$（$0<\rho\leq\frac{1}{10}\rho_1$）使
\begin{equation}\label{eq:3.12}
  \mathcal{H}^{1}\bigl((E\setminus F)\cap B_r(y)\bigr)<2r\cdot 10^{-3}\alpha\qquad(0<r\leq 3\rho)
\end{equation}
且
\begin{equation}\label{eq:3.13}
  \mathcal{H}^{1}\bigl(F\cap B_\rho(y)\bigr)>\frac{1}{4}\cdot 2\rho=\frac{1}{2}\rho.
\end{equation}
必要时缩小 $\rho$，还可设圆盘 $B_\rho(y)$ 的圆周 $\Gamma$ 含有 $F$ 的某个点.

令 $\mathcal{C}$ 为闭圆盘族
\begin{equation}\label{eq:3.14}
  \mathcal{C}=\bigl\{B_r(x):x\in F\cap B_\rho(y),\ 0<r<2\rho
  \ \text{且}\ \mathcal{H}^{1}\bigl((E\setminus F)\cap B_r(x)\bigr)\geq\alpha r\bigr\}.
\end{equation}
由\cref{lem:1.9}，可在 $\mathcal{C}$ 中找到有限或可数的不交闭圆盘族 $\{B_i\}$，使 $\bigcup_{B\in\mathcal{C}}B\subset\bigcup_i B_i'$，其中 $B_i'$ 是与 $B_i$ 同心、半径为其 $5$ 倍的闭圆盘.此外，可使圆盘 $\{B_i'\}$ 半不交.记 $\Gamma_i$ 为 $B_i'$ 的圆周.令
\[
  G=\bigl(F\cap B_\rho(y)\bigr)\cup\Gamma\cup\bigcup_{i}B_i',
\]
并定义
\begin{equation}\label{eq:3.15}
  H=\Bigl(G\setminus\bigcup_{i}B_i'\Bigr)\cup\bigcup_{i}\Gamma_i,
\end{equation}
见图~\ref{fig:3.2}.证明的其余部分分为六步，验证 $H$ 是具有所述性质的连续统.

(a) $G$ 是闭的：由于 $F\cap B_\rho(y)$ 与 $\Gamma$ 闭，只需证 $\overline{\bigcup_{i}B_i'}\subset G$.若 $z\in\overline{\bigcup_{i}B_i'}$，则要么 $z\in\bigcup_{i=1}^{k}B_i'\subset G$ 对某个有限的 $k$ 成立，要么 $z$ 是圆盘子列 $\{B_i'\}_i$ 中点的极限.把面积求和可知 $\sum_i|B_i|^2<\infty$，故当 $i\to\infty$ 时 $|B_i'|\to 0$；于是 $z$ 是这些圆盘的圆心的极限点，而这些圆心都位于闭集 $F\cap B_\rho(y)$ 之中.

(b) $G$ 连通：设 $G=G_1\cup G_2$，其中 $G_1,G_2$ 是非空不交闭集，并由此导出矛盾.$G$ 的任何连通子集必含于 $G_1$ 或 $G_2$ 之一，故设 $\Gamma\subset G_1$.每个 $B_i'$ 的圆心都在 $B_\rho(y)$ 内，故若 $B_i'$ 含有 $\Gamma$ 上或 $\Gamma$ 外的点，它必与 $\Gamma$ 相交，从而含于 $G_1$.于是 $G_2\subset\operatorname{int}B_\rho(y)$.

\begin{figure}[htbp]
  \centering
  \includegraphics[width=0.6\textwidth]{fig3-2.png}
  \caption{引理 3.22 证明中的集合 $G$ 与 $H$}
  \label{fig:3.2}
\end{figure}

设 $G_1'$ 为把 $\mathbb{R}^2\setminus B_\rho(y)$ 并入 $G_1$ 后所得的集合，则 $G_1'$ 闭、连通且与 $G_2$ 不交.由于 $G_1'$ 与 $G_2$ 都含 $F$ 的点，设 $x_1\in F\cap G_1'$ 与 $x_2\in F\cap G_2$ 为使距离 $r=|x_1-x_2|$ 最小的点.该下确界可以取得且为正.由于 $\Gamma\subset G_1'$ 含 $F$ 的点，$0<r<2\rho$.公共区域 $R(x_1,x_2)$ 与 $F$ 不交，否则 $r$ 可进一步缩小，故由引理的假设
\[
  0<\alpha r
  \leq\mathcal{H}^{1}\bigl(E\cap R(x_1,x_2)\bigr)
  =\mathcal{H}^{1}\bigl((E\setminus F)\cap R(x_1,x_2)\bigr)
  \leq\mathcal{H}^{1}\bigl((E\setminus F)\cap B_r(x_2)\bigr).
\]
由于 $x_2\in F\cap B_\rho(y)$ 且 $r<2\rho$，$B_r(x_2)\in\mathcal{C}$，故 $B_r(x_2)\subset\bigcup_i B_i'\subset G$.于是
\[
  B_r(x_2)=(B_r(x_2)\cap G_1')\cup(B_r(x_2)\cap G_2)
\]
是 $B_r(x_2)$ 的由不交闭集构成的分解，$x_1$ 属于第一个，$x_2$ 属于第二个.显然荒谬，因为 $B_r(x_2)$ 连通.

(c) $H$ 是连续统：由 (a) 与 (b)，$G$ 是连续统，故由\cref{lem:3.21}，$H$ 也是连续统.

(d) $\sum_i|B_i'|\leq\frac{1}{8}\rho$：利用 \eqref{eq:3.14}、$B_i$ 的不交性与 \eqref{eq:3.12}，
\[
  \sum_{i}|B_i'|
  =5\sum_{i}|B_i|
  \leq\frac{10}{\alpha}\sum_{i}\mathcal{H}^{1}\bigl((E\setminus F)\cap B_i\bigr)
  \leq\frac{10}{\alpha}\mathcal{H}^{1}\bigl((E\setminus F)\cap B_{3\rho}(y)\bigr)
  \leq\frac{10}{\alpha}\cdot 6\rho\,10^{-3}\alpha
  \leq\frac{1}{8}\rho.
\]

(e) $\mathcal{H}^1(H)<\infty$：由于 $H\subset E\cup\Gamma\cup\bigcup_i\Gamma_i$，利用 (d) 得
\[
  \mathcal{H}^{1}(H)\leq\mathcal{H}^{1}(E)+2\pi\rho+\pi\sum_{i}|B_i'|<\infty.
\]

(f) $\mathcal{H}^1(H\cap E)>0$：由 $H$ 的定义，
\begin{align*}
  \mathcal{H}^{1}(H\cap E)
  &\geq\mathcal{H}^{1}(G\cap E)-\mathcal{H}^{1}\Bigl(E\cap\bigcup_{i}B_i'\Bigr)
  \geq\mathcal{H}^{1}\bigl(F\cap B_\rho(y)\bigr)-\sum_{i}\mathcal{H}^{1}\bigl(E\cap B_i'\bigr)\\
  &\geq\mathcal{H}^{1}\bigl(F\cap B_\rho(y)\bigr)-2\sum_{i}|B_i'|
  \geq\frac{1}{2}\rho-\frac{1}{4}\rho=\frac{1}{4}\rho>0,
\end{align*}
其中用到 \eqref{eq:3.11}（注意 $|B_i'|<20\rho\leq 2\rho_1$），而最后的估计由 \eqref{eq:3.13} 与 (d) 提供.
\end{proof}

关于 $Z$-集下密度的定理现在是圆对几何的一个简单推论.我们利用如下事实：以下密度与法向凸密度都很高的点为心的圆对，必在其公共区域中包含一个正测度子集.

\begin{theorem}\label{thm:3.23}
设 $E$ 是 $\mathbb{R}^2$ 中的 $Z$-集.则在几乎所有的 $x\in E$ 处 $\underline{D}^1(E,x)\leq\frac{3}{4}$.
\end{theorem}

\begin{proof}
设对某个 $\alpha>0$，集合 $E_1=\{x:D^1(E,x)>\frac{3}{4}+\alpha\}$ 具有正测度.于是利用\cref{thm:2.3}，按通常方式可找到紧 $1$-集 $E_2\subset E_1$ 与 $\rho>0$ 使
\begin{equation}\label{eq:3.16}
  \mathcal{H}^{1}\bigl(E\cap B_r(x)\bigr)>\bigl(\tfrac{3}{4}+\alpha\bigr)2r
  \qquad(x\in E_2,\ 0<r\leq\rho)
\end{equation}
且
\begin{equation}\label{eq:3.17}
  \mathcal{H}^{1}(E\cap U)<(1+\alpha)|U|
  \qquad(x\in E_2\cap U,\ 0<|U|\leq 3\rho).
\end{equation}
设 $E_0$ 是 $E_2$ 的紧子集，满足 $0<\mathcal{H}^1(E_0)<\infty$ 且 $|E_0|\leq\rho$.若 $x_1,x_2\in E_0$，则 $r=|x_1-x_2|\leq\rho$，于是注意到任何圆与 $E$ 的交测度为零，并利用 \eqref{eq:3.16}、\eqref{eq:3.17}（取 $U=B_r(x_1)\cap B_r(x_2)$）得
\begin{align*}
  \mathcal{H}^{1}\bigl(E\cap R(x_1,x_2)\bigr)
  &=\mathcal{H}^{1}\bigl(E\cap B_r(x_1)\bigr)+\mathcal{H}^{1}\bigl(E\cap B_r(x_2)\bigr)
  -\mathcal{H}^{1}\bigl(E\cap(B_r(x_1)\cup B_r(x_2))\bigr)\\
  &\geq 2\bigl(\tfrac{3}{4}+\alpha\bigr)2r-(1+\alpha)3r=\alpha r=\alpha|x_1-x_2|.
\end{align*}
于是对 $E_0\subset E$ 应用\cref{lem:3.22}，推得存在连续统 $H$ 使 $0<\mathcal{H}^1(H\cap E)\leq\mathcal{H}^1(H)<\infty$.由于 $E$ 是 $Z$-集，由\cref{lem:3.20} 这不可能.故当 $\alpha>0$ 时 $\mathcal{H}^1(E_1)=0$，如所需.
\end{proof}

\begin{corollary}\label{cor:3.24}
设 $E$ 是 $\mathbb{R}^2$ 中的\newterm{irregular-1-set}{不规则 $1$-集}{irregular 1-set}.则在几乎所有的 $x\in E$ 处 $\underline{D}^1(E,x)\leq\frac{3}{4}$.
\end{corollary}

\begin{proof}
由\cref{cor:3.9}，$E$ 的任何正测度子集都不含于任何可求长曲线.故 $E$ 是 $Z$-集，结论由\cref{thm:3.23} 立即得到.
\end{proof}

我们已经证明：就下密度而言，正则与不规则 $1$-集的性质截然不同——正则集的下密度为 $1$，而不规则集的下密度几乎处处至多为 $\frac{3}{4}$.不可能构造出在正测度集上下密度严格介于 $\frac{3}{4}$ 与 $1$ 之间的 $1$-集.

$\underline{D}^1(E,x)$ 的精确的本质上界仍未确定.虽然 $\frac{3}{4}$ 是迄今得到的最佳值，但人们普遍相信正确答案是 $\frac{1}{2}$；确实可以构造出使 $D^1(E,x)=\frac{1}{2}$ 几乎处处成立的集合.有趣的是，\textcite{besicovitch1928a} 最初给出的界是 $1-10^{-2576}$，其后在 1938 年的文章中才把它降到 $\frac{3}{4}$.

\textcite{besicovitch1938} 还发展了他的圆对方法以证明不规则集的密度几乎处处不存在，另见 \textcites{besicovitchwalker1931,gillis1934b,morserandolph1944}.

现在可以刻画正则与不规则 $1$-集了.

\begin{theorem}\label{thm:3.25}
不规则 $1$-集是 $Z$-集.\newterm{regular-1-set}{正则 $1$-集}{regular 1-set}由一个 $Y$-集再加上一个 $\mathcal{H}^1$-测度为零的集合组成.
\end{theorem}

\begin{proof}
若 $E$ 不规则，则 $E$ 与每条可求长曲线 $\Gamma$ 的交测度为零，否则 $E\cap\Gamma$ 将是一个 $Y$-集，从而是正则的.

若 $E$ 正则，则 $D^1(E,x)=1$ 几乎处处成立，于是\cref{thm:3.23} 结合\cref{cor:2.6} 蕴含：$E$ 的任何正测度可测子集与某条可求长曲线的交具有正测度.我们用这一事实定义可求长曲线序列 $\{\Gamma_j\}$.选取 $\Gamma_1$ 使
\[
  \mathcal{H}^{1}(\Gamma_1\cap E)\geq\frac{1}{2}\sup\bigl\{\mathcal{H}^{1}(\Gamma\cap E):\Gamma\ \text{可求长}\bigr\}.
\]
设 $\Gamma_1,\ldots,\Gamma_k$ 已经选出且 $E_k=E\setminus\bigcup_{1}^{k}\Gamma_j$ 具有正测度，则令 $\Gamma_{k+1}$ 为满足
\begin{equation}\label{eq:3.18}
  \mathcal{H}^{1}(\Gamma_{k+1}\cap E_k)\geq\frac{1}{2}\sup\bigl\{\mathcal{H}^{1}(\Gamma\cap E_k):\Gamma\ \text{可求长}\bigr\}>0
\end{equation}
的可求长曲线.

只有当对某个 $k$，曲线 $\{\Gamma_j\}_{1}^{k}$ 覆盖了 $E$ 的几乎全部时过程才终止，此时结论显然.否则，
\[
  \infty>\mathcal{H}^{1}(E)\geq\sum_{k}\mathcal{H}^{1}(\Gamma_{k+1}\cap E_k),
\]
故 $\mathcal{H}^1(\Gamma_{k+1}\cap E_k)\to 0$.若 $\mathcal{H}^1\bigl(E\setminus\bigcup_{1}^{\infty}\Gamma_j\bigr)>0$，则可找到可求长曲线 $\Gamma$ 使 $\mathcal{H}^1\bigl(\Gamma\cap(E\setminus\bigcup_{1}^{\infty}\Gamma_j)\bigr)=d>0$.但对某个 $k$ 有 $\mathcal{H}^1(\Gamma_{k+1}\cap E_k)<\frac{1}{2}d$，而按 \eqref{eq:3.18}，$\Gamma$ 本应被优先选为 $\Gamma_{k+1}$.故 $\mathcal{H}^1\bigl(E\setminus\bigcup_{1}^{\infty}\Gamma_j\bigr)=0$，于是 $E$ 由 $Y$-集 $\bigcup_{j}(E\cap\Gamma_j)$ 与一个零测集组成.
\end{proof}

不规则 $1$-集作为 $Z$-集，与每条可求长曲线的交测度为零.然而 \textcite{besicovitch1928a}（第 43 节）表明，不规则集可以嵌入非可求长曲线的可数并之中.

本节最后指出不规则集的``破碎''性质.回忆：若集合的任何两点都不位于同一连通分支中，则称该集合是完全不连通的.于是，给定集合中的任意一对点，存在由两个不交闭子集构成的分解，每个子集含其中一点.

\begin{corollary}\label{cor:3.26}
不规则 $1$-集 $E$ 是\newterm{totally-disconnected}{完全不连通的}{totally disconnected}.
\end{corollary}

\begin{proof}
若 $x$ 与 $y$ 同属于 $E$ 的某个闭连通子集 $F$，则由\cref{lem:3.4} 有 $|x-y|\leq\mathcal{H}^1(F)<\infty$.但 $E$ 是 $Z$-集，故由\cref{lem:3.20} 有 $\mathcal{H}^1(F)=\mathcal{H}^1(E\cap F)=0$.于是实际上 $x=y$.
\end{proof}

\section{切线性质}

与密度性质一样，正则集与不规则集的切线性质截然不同.回忆切线的定义 \eqref{eq:3.4}，正则集的结果已经清楚：

\begin{theorem}\label{thm:3.27}
$\mathbb{R}^n$ 中的正则 $1$-集 $E$ 在几乎所有的点处有切线.
\end{theorem}

\begin{proof}
由\cref{thm:3.25}，$E$ 由一个 $Y$-集与一个零测集组成，而由\cref{cor:3.10}，任何 $Y$-集在几乎所有点处有切线.
\end{proof}

为研究平面中不规则集的切线性质，我们使用 \textcite{besicovitch1938}（第 10 节）的一个精巧论证.首先需要一个几何引理.

\begin{lemma}\label{lem:3.28}
设 $\theta$ 是 $\mathbb{R}^2$ 中与直线 $L$ 垂直的单位向量.设 $P$ 是边与方向 $\pm\theta$ 成角 $\phi$ 的平行四边形，$y$ 与 $z$ 是 $P$ 的一对相对顶点，如图~\ref{fig:3.3}.则 $|y-z|\leq d/\sin\phi$，其中 $d$ 是 $P$ 到 $L$ 上投影的长度.
\end{lemma}

\begin{figure}[htbp]
  \centering
  \includegraphics[width=0.4\textwidth]{fig3-3.png}
  \caption{引理 3.28 中的平行四边形}
  \label{fig:3.3}
\end{figure}

\begin{proof}
设 $w$ 为 $P$ 的第三个顶点.则 $|y-z|\leq|y-w|+|w-z|$，且 $d=(|y-w|+|w-z|)\sin\phi$，由此得结论.
\end{proof}

下一条定理本质上是 Besicovitch 关于不规则集的切线结果.证明的想法是把不规则点集约化为一条可求长曲线的子集，而后者必有零测度.

\begin{theorem}\label{thm:3.29}
设 $E$ 是 $\mathbb{R}^2$ 中的不规则 $1$-集.则给定 $\theta$ 与 $0<\phi<\pi/2$，
\begin{equation}\label{eq:3.19}
  \bar{D}^{1}(E,x,\theta,\phi)+\bar{D}^{1}(E,x,-\theta,\phi)\geq\frac{1}{6}\sin\phi
\end{equation}
对几乎所有的 $x\in E$ 成立.
\end{theorem}

\begin{proof}
取 $\rho,\delta_+,\delta_->0$，并令 $F_0=F_0(\delta_+,\delta_-,\rho)$ 为 $E$ 中使
\begin{equation}\label{eq:3.20}
  \left.
  \begin{aligned}
    \mathcal{H}^{1}\bigl(E\cap S_r(x,\theta,\phi)\bigr)&\leq 2r\delta_+\\
    \mathcal{H}^{1}\bigl(E\cap S_r(x,-\theta,\phi)\bigr)&\leq 2r\delta_-
  \end{aligned}
  \right\}
\end{equation}
对一切 $r\leq\rho$ 成立的 $x$ 之集.照例 $F_0$ 可测；我们将证明若 $F_0$ 有正测度，则 $\delta_++\delta_-$ 不能太小.

若 $\mathcal{H}^1(F_0)>0$，则由 $\mathcal{H}^1$ 的正则性可找到 $F_0$ 的正测度紧子集 $F_1$.进一步，由\cref{thm:2.3} 或直接由 $\mathcal{H}^1$ 的定义，给定 $\eta>0$，可找到闭凸集 $U$ 使 $0<|U|\leq\rho$ 且
\begin{equation}\label{eq:3.21}
  \mathcal{H}^{1}(F)>(1-\eta)|U|,
\end{equation}
其中 $F=F_1\cap U$.从此我们只在 $U$ 内工作.

由于 $F$ 闭，可选取 $y_1$ 与 $z_1$ 为 $F$ 中连接线段与 $\theta$ 的夹角不超过 $\phi$ 的最远点对之一，使得
\[
  r_1=|y_1-z_1|
  =\sup\bigl\{|y-z|:z\in F\cap S(y,\theta,\phi)\ \text{且}\ y\in F\cap S(z,-\theta,\phi)\bigr\},
\]
其中 $S(y,\theta,\phi)$ 表示无穷扇形.$r_1$ 的极大性保证
\[
  F\cap S(y_1,\theta,\phi)=F\cap S_{r_1}(y_1,\theta,\phi)
\]
且
\[
  F\cap S(z_1,-\theta,\phi)=F\cap S_{r_1}(z_1,-\theta,\phi).
\]
由于 $r_1\leq\rho$ 且 $F\subset E$，由 \eqref{eq:3.20} 得
\begin{equation}\label{eq:3.22}
  \left.
  \begin{aligned}
    \mathcal{H}^{1}\bigl(F\cap S(y_1,\theta,\phi)\bigr)&\leq 2r_1\delta_+\\
    \mathcal{H}^{1}\bigl(F\cap S(z_1,-\theta,\phi)\bigr)&\leq 2r_1\delta_-.
  \end{aligned}
  \right\}
\end{equation}
令 $P_1$ 为闭平行四边形
\[
  P_1=S(y_1,\theta,\phi)\cap S(z_1,-\theta,\phi)
\]
并令 $Q_1$ 为开区域
\[
  Q_1=\operatorname{int}\bigl\{S(y_1,\theta,\phi)\cup S(z_1,-\theta,\phi)\bigr\},
\]
见图~\ref{fig:3.4}.由 \eqref{eq:3.22},
\[
  \mathcal{H}^{1}(F\cap Q_1)\leq 2r_1(\delta_++\delta_-).
\]

\begin{figure}[htbp]
  \centering
  \includegraphics[width=0.42\textwidth]{fig3-4.png}
  \caption{\cref{thm:3.29} 证明中的扇形与区域}
  \label{fig:3.4}
\end{figure}

集合 $U\setminus Q_1$ 紧且至多有两个分支，因此只要它在其同一分支中含有连接线段与 $\theta$ 的夹角至多为 $\phi$ 的一对 $F$ 中的点，我们就可用同样的方法得到位于 $Q_1$ 之外的平行四边形 $P_2$ 与开区域 $Q_2$.于是对 $j=1,2,\ldots$，我们在 $U\setminus\bigcup_{i=1}^{j-1}Q_i$ 的同一分支中找到 $F$ 的点 $y_j$ 与 $z_j$，使 $r_j=|y_j-z_j|$ 尽可能大，且
\begin{equation}\label{eq:3.23}
  \mathcal{H}^{1}(F\cap Q_j)\leq 2r_j(\delta_++\delta_-),
\end{equation}
其中
\[
  Q_j=\operatorname{int}\bigl\{S(y_j,\theta,\phi)\cup S(z_j,-\theta,\phi)\bigr\},
  \qquad
  P_j=S(y_j,\theta,\phi)\cap S(z_j,-\theta,\phi).
\]
这一过程或者无限进行下去，或者直到没有合适的点对为止.

设 $L$ 为与 $\theta$ 垂直的直线.若 $j>i$，则 $P_j$ 与 $Q_i$ 不交，故平行四边形 $\{P_j\}$ 在 $L$ 上的投影是内部不交的线段.用\cref{lem:3.28} 容易验证
\begin{equation}\label{eq:3.24}
  \sum_{j}r_j=\sum_{j}|y_j-z_j|\leq\frac{3|U|}{\sin\phi};
\end{equation}
因子 `3' 顾及了至多两个 $P_j$ 在 $U$ 两侧重叠的可能性.

设 $y$ 与 $z$ 为 $F\setminus\bigcup_{i}Q_i$ 中两个不同的点.若 $y$ 与 $z$ 位于某个 $Q_j$ 的两侧，则 $[y,z]$ 与 $\theta$ 的夹角至少为 $\phi$.若不然，且 $[y,z]$ 与 $\theta$ 的夹角小于 $\phi$，则 $y$ 与 $z$ 在每个阶段都本是 $y_j$ 与 $z_j$ 的候选者，故选取过程不会终止，且对一切 $j$ 有 $|y-z|\leq|y_j-z_j|$.这是不可能的，因为 \eqref{eq:3.24} 蕴含 $|y_j-z_j|\to 0$.于是我们得出：连接 $F\setminus\bigcup_{i}Q_i$ 中任意两点对线段与 $L$ 的夹角至多为 $\frac{1}{2}\pi-\phi$.由\cref{lem:3.11}, $F\setminus\bigcup_{i}Q_i$ 含于一条可求长曲线之中.由于 $F\setminus\bigcup_{i}Q_i$ 是 $E$ 的子集，而由\cref{thm:3.25} $E$ 是 $Z$-集，故
\[
  \mathcal{H}^{1}\Bigl(F\setminus\bigcup_{i}Q_i\Bigr)=0.
\]
于是利用 \eqref{eq:3.23}、\eqref{eq:3.24} 与 \eqref{eq:3.21},
\[
  \mathcal{H}^{1}(F)
  \leq\sum_{j}\mathcal{H}^{1}(F\cap Q_j)
  \leq 2(\delta_++\delta_-)\sum_{j}r_j
  \leq\frac{6(\delta_++\delta_-)|U|}{\sin\phi}
  \leq\frac{6(\delta_++\delta_-)}{(1-\eta)\sin\phi}\,\mathcal{H}^{1}(F).
\]
对任意 $\eta>0$ 都存在使上式成立的满足 $\mathcal{H}^1(F)>0$ 的 $F$，故
\[
  (\delta_++\delta_-)\geq\frac{1}{6}\sin\phi.
\]
于是若 $\delta_++\delta_-<\frac{1}{6}\sin\phi$，则由 \eqref{eq:3.20} 定义的集合 $F_0(\delta_+,\delta_-,\rho)$ 必测度为零，证明完成.
\end{proof}

\begin{corollary}\label{cor:3.30}
设 $E$ 是 $\mathbb{R}^2$ 中的不规则 $1$-集.则对几乎所有的 $x\in E$,
\begin{equation}\label{eq:3.25}
  \bar{D}^{1}(E,x,\theta,\phi)+\bar{D}^{1}(E,x,-\theta,\phi)\geq\frac{1}{6}\sin\phi>\frac{1}{10}\phi
\end{equation}
对一切 $\theta$ 与一切 $0<\phi<\frac{1}{2}\pi$ 成立.
\end{corollary}

\begin{proof}
由\cref{thm:3.29}, \eqref{eq:3.19} 对 $\{(\theta,\phi)\}$ 的某个可数稠密子集对几乎所有的 $x$ 成立.于是通过逼近，对几乎所有的 $x$，\eqref{eq:3.19} 对一切 $\theta$ 与 $\phi$ 成立.
\end{proof}

\eqref{eq:3.19} 或 \eqref{eq:3.25} 右端的系数当然可以改进.事实上，Besicovitch 的证明（用了略优于\cref{lem:3.28} 与 \eqref{eq:3.24} 的估计）给出 $\frac{1}{3}\sin\phi$，而且当 $\phi\geq\frac{1}{3}\pi$ 时给出 $\frac{1}{2}\sin\phi$.Besicovitch 猜测最佳系数可能大至 $2^{-\frac{1}{3}}$（Besicovitch (1982a, 第 11 节) 中的一个例子表明它不能比这更大）.\footnote{原书参考文献表中并无 Besicovitch (1982a) 这一条目，疑为 \citeyear{besicovitch1928a} 之误印，请读者鉴察.——译者注.}然而，任何正值都足以导出主要的切线结果.

\begin{corollary}\label{cor:3.31}
在平面不规则 $1$-集的几乎所有点处不存在切线.
\end{corollary}

\begin{proof}
由\cref{cor:3.30} 与切线的定义 \eqref{eq:3.4} 立得.
\end{proof}

如我们所见，正则 $1$-集可以作为可求长曲线的子集得到.不规则集的例子也许不那么显然，因此我们现在给出一个不规则 $1$-集的简单构造.

\begin{theorem}\label{thm:3.32}
$\mathbb{R}^2$ 中存在不规则 $1$-集.
\end{theorem}

\begin{proof}
设 $f$ 是 $[0,1)$ 上由
\[
  f(.x_1x_2\ldots)=.y_1y_2\ldots
\]
定义的实函数，其中 $x_1x_2\ldots$ 与 $y_1y_2\ldots$ 以 $4$ 为底书写（约定 $x_i$ 不最终全等于 $3$），使得 $y_i=5-x_i\pmod{4}$.（见图~\ref{fig:3.5}.）令
\[
  E=\{(x,f(x)):0\leq x<1\}.
\]
$E$ 是 Borel 集，且由于 $E$ 到 $x$ 轴的投影长度为 $1$，得 $1\leq\mathcal{H}^1(E)$.另一方面，对每个 $k$，$E$ 可用 $4^k$ 个边长为 $4^{-k}$ 的正方形覆盖，故 $\mathcal{H}^1(E)\leq\sqrt{2}$.于是 $E$ 是一个 $1$-集.容易看出任何可求长曲线（它在几乎处处有切线）与 $E$ 的交长度为零.故 $E$ 是 $Z$-集，是不规则的.另外，不规则性也可由 $E$ 在与 $x$ 轴成 $\pm 45^{\circ}$ 的两个方向上投影成勒贝格测度为零的集合这一事实得到；见\cref{cor:6.14}.
\end{proof}

\begin{figure}[htbp]
  \centering
  \includegraphics[width=0.42\textwidth]{fig3-5.png}
  \caption{\cref{thm:3.32} 中函数 $y=f(x)$ 的图像}
  \label{fig:3.5}
\end{figure}

\section{高维集合}

把该理论推广到 $\mathbb{R}^n$（$n\geq 3$）的整数维数子集，直到 Besicovitch 的第一篇基础性论文 \citeyear{besicovitch1928a} 发表约 47 年之后才完成.这是因为必须克服相当大的技术困难；例如，``有限 $\mathcal{H}^s$-测度的连续统是可求长曲线之并''这一定理当 $s$ 为大于 $1$ 的整数时没有类似物.

称 $\mathbb{R}^n$ 中的 $s$-集为\newterm{countably rectifiable}{可数可求长的}{countably rectifiable}，如果它形如
\[
  \bigcup_{j=1}^{\infty}f_j(E_j)\cup G,
\]
其中 $\mathcal{H}^s(G)=0$，且每个 $f_j$ 都是从 $\mathbb{R}^s$ 的有界子集 $E_j$ 到 $\mathbb{R}^n$ 的 \newterm{lipschitz-mapping}{Lipschitz 函数}{Lipschitz mapping}.（Lipschitz 函数 $f$ 要求对某个常数 $c$ 有 $|f(t)-f(u)|\leq c|t-u|$.）这是 Besicovitch 的``$Y$-集加上一个零测集''的直接类似物.

$\mathbb{R}^n$ 中的 $s$-集 $E$ 称为在 $x$ 处有切线，如果 $\bar{D}^s(E,x)>0$ 且存在过 $x$ 的 $s$ 维平面 $\Pi$，使对一切 $\phi>0$,
\begin{equation}\label{eq:3.26}
  \lim_{r\to 0}r^{-s}\mathcal{H}^{s}\bigl(E\cap(B_r(x)\setminus S(x,\Pi,\phi))\bigr)=0,
\end{equation}
其中 $S(x,\Pi,\phi)$ 表示使 $[y,x]$ 与 $\Pi$ 的夹角至多为 $\phi$ 的点 $y\in\mathbb{R}^n$ 之集.若 $\underline{D}^s(E,x)>0$ 且把 \eqref{eq:3.26} 中的 `lim' 换成 `$\varliminf$'，则称 $E$ 在 $x$ 处有\newterm{weak tangent}{弱切线}{weak tangent}.

\begin{theorem}\label{thm:3.33}
设 $E$ 是 $\mathbb{R}^n$ 中的 $s$-集，$s$ 为整数.则下述命题等价：

(a) $E$ 是正则的.

(b) $E$ 是可数可求长的.

(c) $E$ 在几乎所有的点处有切线.
\end{theorem}

如我们所见，Besicovitch（\citeyear{besicovitch1928a,besicovitch1938}） 对 $\mathbb{R}^2$ 中的 $1$-集证明了这一点.\textcite{federer1947} 在一篇极其复杂的文章中证明了 (b) 与 (c) 的等价性，而由这一点及切线的定义可得 (b) 蕴含 (a).（Federer 的 restrictedness 概念等价于切线的存在性.）\textcite{moore1950} 完成了 $\mathbb{R}^n$ 中 $1$-集的证明，\textcite{marstrand1961} 完成了 $\mathbb{R}^3$ 中 $2$-集的证明，而 \textcite{mattila1975a} 最终完成了 $\mathbb{R}^n$ 中 $s$-集的证明.对更一般的测度 $\mu$ 也得到了此类性质的结果.这里必须把正则性单纯定义为密度 $\lim_{r\to 0}r^{-s}\mu(B_r(x))$ $\mu$-几乎处处存在，而不涉及极限的实际值.上文提到的 Federer 与 Moore 的文章考虑了这样的推广，\textcite{morserandolph1944} 考虑了 $\mathbb{R}^2$ 中 $s=1$ 的子集，\textcite{marstrand1964} 则考虑了一般情形.特别地，Marstrand 证明对一般测度，正则性蕴含弱切线（按显然的方式定义）几乎处处存在；自然的猜想是 `弱'一词可以省去.

这些证明中的许多都与\cref{chap:6} 中讨论的投影性质的高维类似物紧密相关.

当 $s$ 是整数且 $1<s<n$ 时，把曲线理论推广到 $\mathbb{R}^n$ 中的 $s$ 维\newterm{surface}{曲面}{surface}存在相当的困难.这一理论远比曲线理论复杂；事实上，对``曲面''与``面积''等术语所赋予的精确含义需要格外小心.

$\mathbb{R}^3$ 中的曲面已被研究多年，很长时间里公认的曲面面积定义是 \newterm{lebesgue-frechet-area}{Lebesgue--Fréchet 定义}{Lebesgue-Fréchet area}，即作为内接多面体面积的上确界，与曲线长度的定义相类比.正是 \textcite{besicovitch1945b} 证明了这一定义是全然不适当的，它甚至不具备测度的基本可加性质.他提出最令人满意的面积定义是 Hausdorff 测度 $\mathcal{H}^2$，并着手由此出发处理曲面面积的各类问题.

由此产生的一个问题是如何把 Gołąb 半连续性定理（\cref{thm:3.18}）推广到曲面.显然没有直接的类似物，因为任何曲面都可以在 Hausdorff 度量下用一列可求长曲线逼近，而后者理应有零面积.\textcite{besicovitch1949} 证明了：若 $\{E_j\}$ 是收敛于参数曲面 $E$ 的参数曲面序列，且极限曲面 $E$ 满足合理的光滑性条件，则 $\mathcal{H}^2(E)\leq\varliminf_{j\to\infty}\mathcal{H}^2(E_j)$.把 Gołąb 定理向高维集合的一个更近的推广归功于 \textcite{vitushkin1966}，另见 \textcite{ivanov1975}.我们如下定义 $\mathbb{R}^n$ 的紧子集 $E$ 的\newterm{variations}{变差}{variations}.令 $V^0(E)$ 为 $E$ 的连通分支的个数，于是当 $V^0(E)=1$ 时 $E$ 是连续统.对 $1\leq s\leq n$，令 $V^s(E)=c\int V^0(E\cap\Pi)\,\mathrm{d}\Pi$，其中 $\Pi$ 取遍一切 $(n-s)$ 维平面，积分关于所有这类平面的空间上的自然不变测度（出现于积分几何中的那个）进行.规范化常数 $c$ 只依赖于 $n$ 与 $s$，其选取使得当 $E$ 是 $s$ 维方体时 $V^s(E)$ 就是其体积.于是对充分光滑的曲面或多面体曲面，$V^s$ 就是 $s$ 维体积.这些变差起初用于凸集的研究，但也可以卓有成效地用于一般紧集.\newterm{vitushkin-semicontinuity}{Vitushkin 半连续性定理}{semicontinuity theorem, Vitushkin's}断言：若 $\{E_j\}$ 是 $\mathbb{R}^n$ 中在 Hausdorff 度量下收敛于 $E$ 的紧集序列，则只要数 $\{V^i(E_j)\}$ 对 $0\leq i<s$ 一致有界，就有 $V^s(E)\leq\varliminf_{j\to\infty}V^s(E_j)$.当 $s=1$ 时它化为 Gołąb 定理.

一个特别有意义的问题是求\newterm{minimal-surface}{``极小曲面''}{minimal surface}，即张成 $\mathbb{R}^3$ 中一条闭曲线的面积最小的曲面.这需要对``曲面''作仔细的定义，随后问题便适合用 Hausdorff 测度来处理.\newterm{plateau-problem}{``Plateau 问题''}{Plateau problem}——面积最小的曲面是否总存在——最终由 \textcite{reifenberg1960} 以肯定的方式解决.（该文也列出了更早的文献.）\textcite{besicovitch1950} 对其中一些专题作了综述.

若 $E$ 是 $\mathbb{R}^n$ 中满足 $\mathcal{L}^n(E)>0$ 的有界勒贝格可测子集，$E$ 的\newterm{density-boundary}{密度边界}{density boundary} $\partial E$ 由 $\mathbb{R}^n$ 中使 $E$ 的勒贝格密度或者不存在、或者取 $0$ 与 $1$ 以外的值的那些点组成.勒贝格密度定理（\cref{thm:1.13}）蕴含 $\mathcal{L}^n(\partial E)=0$.然而在相反的方向上，当 $E$ 的边界高度不规则时，可以证明 $\mathcal{H}^{n-1}(\partial E)>0$ 这一非平凡结果.

人们还设计了多种其他曲面测度.除已提到者外，常用的还包括积分几何测度与 \newterm{de-gorgi-perimeter}{de Giorgi 周长}{de Gorgi perimeter}；见 \textcite{federer1969}，那里详细定义并考察了这类测度.

\begin{exercises}
\item 设 $\psi:[a,b]\to\mathbb{R}^n$ 是（不必连续的）有界变差函数，即
$\operatorname{var}[a,b]=\sup\sum_{i=1}^{m}|\psi(t_i)-\psi(t_{i-1})|<\infty$，
其中上确界取遍 $[a,b]$ 的一切分割 $a=t_0<t_1<\cdots<t_m=b$.证明 $\psi[a,b]$ 含于一条（可能自交的）可求长曲线.（提示：用 $\operatorname{var}[a,t]$ 定义一个从 $[0,\operatorname{var}[a,b]]$ 到 $\mathbb{R}^n$ 的连续函数.）进而证明 $\mathcal{H}^1(\psi[a,b])<\infty$.
\item 证明任何弧连通集都是连通的.给出 $\mathbb{R}^2$ 中一个不是弧连通的连续统的例子.（比较\cref{lem:3.12}.）
\item 验证 Hausdorff 度量是一个度量.
\item 给出 $\mathbb{R}^2$ 中在 Hausdorff 度量下收敛于 $E$ 的连续统序列 $\{E_j\}$ 的例子，使 $\mathcal{H}^1(E_j)\to\infty$ 而 $\mathcal{H}^1(E)<\infty$.
\item 设 $E$ 是满足 $\mathcal{H}^1(E)<\infty$ 的连续统.证明存在树 $F\subset E$ 使 $\mathcal{H}^1(E\setminus F)$ 任意小.进而证明 $E$ 含于一条长度至多为 $2\mathcal{H}^1(E)$ 的（可能自交的）曲线之中.
\item 设 $E$ 是不规则 $1$-集.给定 $\varepsilon,c>0$，证明存在 $\delta_0>0$，使得对任何满足 $\mathcal{H}^1(F)<c$ 的连续统 $F$，当 $0<\delta<\delta_0$ 时有 $\mathcal{H}^1\bigl(E\cap[F]_{\delta}\bigr)<\varepsilon$，其中 $[F]_{\delta}$ 是 $F$ 的 $\delta$-平行体.
\item 证明存在 Hausdorff 维数为 $2$ 的紧集 $F\subset\mathbb{R}^2$，使得含于 $F$ 中的每个 $1$-集都是不规则的.
\item 在 $\mathbb{R}^2$ 中构造一个完全不连通的正则 $1$-集.（比较\cref{cor:3.26}.）
\end{exercises}
